% GENERATED FILE. Edit canonical lessons, then run npm run generate:books.
% canonical-source-hash: fnv1a64:2f69d69c296d7f64
% canonical-lessons: KA-C513-cetarisiko, KA-C513-unisu, KA-C513-badita, KA-C513-araike, KA-C513-kasaya

\chapter{Getting Better}
\label{ch:ka-getting-better}
\hlchaptermodality{513}

\begin{chapteropening}
\textbf{By the end of this chapter:} \emph{I can say to recover, to feed, a heartbeat, nursing and a herbal decoction.}

Complete the last lesson of chapter 513: I can say to recover, to feed, a heartbeat, nursing and a herbal decoction.
\end{chapteropening}

\section[\texorpdfstring{\kn{ಚೇತರಿಸಿಕೊ} (cētarisiko)}{ಚೇತರಿಸಿಕೊ (cētarisiko)}]{\kn{ಚೇತರಿಸಿಕೊ} (cētarisiko) — to recover; to pick up}

\label{lesson:KA-C513-cetarisiko}



\begin{quote}
Before the new one: say the Kannada for a stroke, then the Kannada for to faint.
\end{quote}

\subsection*{You'll want to know: \kn{ಚೇತರಿಸಿಕೊ}}
\textbf{\kn{ಚೇತರಿಸಿಕೊ}} — \emph{cētarisiko} — \textquotedblleft{}to recover; to pick up\textquotedblright{}.

\begin{quote}
\kn{ಆಪರೇಷನ್} \kn{ನಂತರ} \kn{ಅಮ್ಮ} \kn{ಬೇಗ} \kn{ಚೇತರಿಸಿಕೊಂಡರು}.
\end{quote}

\begin{quote}
\emph{āparēṣan nantara amma bēga cētarisikoṇḍaru.}
\end{quote}

\begin{quote}
After the operation mother recovered quickly.
\end{quote}

Past (I …): \kn{ಚೇತರಿಸಿಕೊಂಡೆ}, \emph{cētarisikoṇḍe}.

From what: -\kn{ಇಂದ}. \kn{ಚೇತರಿಕೆ} is recovery.

\begin{cousinweb}
\textbf{\kn{ಚೇತರಿಸಿಕೊ}} (\emph{cētarisiko}) is usually traced to Sanskrit \emph{cetana}, \textquotedblleft{}conscious, sentient; the mind\textquotedblright{} (Monier-Williams), with -\emph{isi-ko}: to bring life back into oneself. The exact path is not settled.
\end{cousinweb}

\begin{grammarlens}[title={What you've built}]
One more word for describing what you lived and felt.
\end{grammarlens}

\subsection*{Your turn}
\begin{itemize}
\raggedright
  \item \emph{Say it:} \emph{cētarisiko}
  \item \emph{Say it:} \emph{cētarisiko}, once more
  \item \emph{Say it:} \emph{āparēṣan nantara amma bēga cētarisikoṇḍaru.} — After the operation mother recovered quickly.
  \item \emph{Recall:} say the Kannada for a stroke, then the Kannada for to faint, then say \emph{cētarisiko} again
\end{itemize}

\begin{tcolorbox}[breakable,colback=teal!4,colframe=teal!35!black,title={Before you move on}]
What does \kn{ಚೇತರಿಸಿಕೊ} mean? (\textquotedblleft{}To recover; to pick up\textquotedblright{}.) Say it once more.
\end{tcolorbox}
\section[\texorpdfstring{\kn{ಉಣಿಸು} (uṇisu)}{ಉಣಿಸು (uṇisu)}]{\kn{ಉಣಿಸು} (uṇisu) — to feed (a child or the sick)}

\label{lesson:KA-C513-unisu}



\begin{quote}
Before the new one: say the Kannada for to faint, then the Kannada for to recover.
\end{quote}

\subsection*{You'll want to know: \kn{ಉಣಿಸು}}
\textbf{\kn{ಉಣಿಸು}} — \emph{uṇisu} — \textquotedblleft{}to feed (a child or the sick)\textquotedblright{}.

\begin{quote}
\kn{ಅಮ್ಮ} \kn{ಜ್ವರ} \kn{ಬಂದ} \kn{ಮಗನಿಗೆ} \kn{ಮೆತ್ತಗಿನ} \kn{ಅನ್ನ} \kn{ಉಣಿಸಿದರು}.
\end{quote}

\begin{quote}
\emph{amma jvara banda maganige mettagina anna uṇisidaru.}
\end{quote}

\begin{quote}
Mother fed her feverish son soft rice.
\end{quote}

Past (I …): \kn{ಉಣಿಸಿದೆ}, \emph{uṇiside}.

The one fed takes -\kn{ಗೆ}. It is the causative of the old verb \kn{ಉಣ್}, to eat a meal, as in \kn{ಊಟ}.

\begin{cousinweb}
\textbf{\kn{ಉಣಿಸು}} (\emph{uṇisu}) is native Dravidian: DEDR 600 lists Kannada \emph{uṇ}, \emph{uṇṇu}, \textquotedblleft{}to eat a meal; of children, to take mother's milk\textquotedblright{}; -\emph{isu} makes \textquotedblleft{}to make someone eat\textquotedblright{}.
\end{cousinweb}

\begin{grammarlens}[title={What you've built}]
One more word for describing what you lived and felt.
\end{grammarlens}

\subsection*{Your turn}
\begin{itemize}
\raggedright
  \item \emph{Say it:} \emph{uṇisu}
  \item \emph{Say it:} \emph{uṇisu}, once more
  \item \emph{Say it:} \emph{amma jvara banda maganige mettagina anna uṇisidaru.} — Mother fed her feverish son soft rice.
  \item \emph{Recall:} say the Kannada for to faint, then the Kannada for to recover, then say \emph{uṇisu} again
\end{itemize}

\begin{tcolorbox}[breakable,colback=teal!4,colframe=teal!35!black,title={Before you move on}]
What does \kn{ಉಣಿಸು} mean? (\textquotedblleft{}To feed (a child or the sick)\textquotedblright{}.) Say it once more.
\end{tcolorbox}
\section[\texorpdfstring{\kn{ಬಡಿತ} (baḍita)}{ಬಡಿತ (baḍita)}]{\kn{ಬಡಿತ} (baḍita) — a heartbeat; a throb}

\label{lesson:KA-C513-badita}



\begin{quote}
Before the new one: say the Kannada for to recover, then the Kannada for to feed.
\end{quote}

\subsection*{You'll want to know: \kn{ಬಡಿತ}}
\textbf{\kn{ಬಡಿತ}} — \emph{baḍita} — \textquotedblleft{}a heartbeat; a throb\textquotedblright{}.

\begin{quote}
\kn{ವೈದ್ಯರು} \kn{ಮಗುವಿನ} \kn{ಹೃದಯದ} \kn{ಬಡಿತ} \kn{ಕೇಳಿದರು}.
\end{quote}

\begin{quote}
\emph{vaidyaru maguvina hṛdayada baḍita kēḷidaru.}
\end{quote}

\begin{quote}
The doctor listened to the baby's heartbeat.
\end{quote}

\kn{ಎದೆ} \kn{ಬಡಿತ}, the beating of the heart; fear makes it \kn{ಜೋರು}.

\begin{cousinweb}
\textbf{\kn{ಬಡಿತ}} (\emph{baḍita}) is native Dravidian: DEDR 5224 lists Kannada \emph{baḍi}, \textquotedblleft{}to beat, strike\textquotedblright{}, with nouns for beating; -\emph{ta} makes the noun.
\end{cousinweb}

\begin{grammarlens}[title={What you've built}]
One more word for describing what you lived and felt.
\end{grammarlens}

\subsection*{Your turn}
\begin{itemize}
\raggedright
  \item \emph{Say it:} \emph{baḍita}
  \item \emph{Say it:} \emph{baḍita}, once more
  \item \emph{Say it:} \emph{vaidyaru maguvina hṛdayada baḍita kēḷidaru.} — The doctor listened to the baby's heartbeat.
  \item \emph{Recall:} say the Kannada for to recover, then the Kannada for to feed, then say \emph{baḍita} again
\end{itemize}

\begin{tcolorbox}[breakable,colback=teal!4,colframe=teal!35!black,title={Before you move on}]
What does \kn{ಬಡಿತ} mean? (\textquotedblleft{}A heartbeat; a throb\textquotedblright{}.) Say it once more.
\end{tcolorbox}
\section[\texorpdfstring{\kn{ಆರೈಕೆ} (āraike)}{ಆರೈಕೆ (āraike)}]{\kn{ಆರೈಕೆ} (āraike) — nursing; looking after}

\label{lesson:KA-C513-araike}



\begin{quote}
Before the new one: say the Kannada for to feed, then the Kannada for a heartbeat.
\end{quote}

\subsection*{You'll want to know: \kn{ಆರೈಕೆ}}
\textbf{\kn{ಆರೈಕೆ}} — \emph{āraike} — \textquotedblleft{}nursing; looking after\textquotedblright{}.

\begin{quote}
\kn{ಅಮ್ಮನ} \kn{ಆರೈಕೆಯಿಂದ} \kn{ನಾನು} \kn{ಬೇಗ} \kn{ಗುಣವಾದೆ}.
\end{quote}

\begin{quote}
\emph{ammana āraikeyinda nānu bēga guṇavāde.}
\end{quote}

\begin{quote}
With mother's nursing, I got well quickly.
\end{quote}

\kn{ಆರೈಕೆ} \kn{ಮಾಡು}, to nurse, to tend.

\begin{cousinweb}
\textbf{\kn{ಆರೈಕೆ}} (\emph{āraike}) is native Dravidian: DEDR 377 lists Kannada \emph{āray}, \emph{ārayu}, \textquotedblleft{}to search, to take care of\textquotedblright{}, with the nouns \emph{ārayike}, \emph{ārayke}.
\end{cousinweb}

\begin{grammarlens}[title={What you've built}]
One more word for describing what you lived and felt.
\end{grammarlens}

\subsection*{Your turn}
\begin{itemize}
\raggedright
  \item \emph{Say it:} \emph{āraike}
  \item \emph{Say it:} \emph{āraike}, once more
  \item \emph{Say it:} \emph{ammana āraikeyinda nānu bēga guṇavāde.} — With mother's nursing, I got well quickly.
  \item \emph{Recall:} say the Kannada for to feed, then the Kannada for a heartbeat, then say \emph{āraike} again
\end{itemize}

\begin{tcolorbox}[breakable,colback=teal!4,colframe=teal!35!black,title={Before you move on}]
What does \kn{ಆರೈಕೆ} mean? (\textquotedblleft{}Nursing; looking after\textquotedblright{}.) Say it once more.
\end{tcolorbox}
\section[\texorpdfstring{\kn{ಕಷಾಯ} (kaṣāya)}{ಕಷಾಯ (kaṣāya)}]{\kn{ಕಷಾಯ} (kaṣāya) — a herbal decoction}

\label{lesson:KA-C513-kasaya}



\begin{quote}
Before the new one: say the Kannada for a heartbeat, then the Kannada for nursing.
\end{quote}

\subsection*{You'll want to know: \kn{ಕಷಾಯ}}
\textbf{\kn{ಕಷಾಯ}} — \emph{kaṣāya} — \textquotedblleft{}a herbal decoction\textquotedblright{}.

\begin{quote}
\kn{ಶೀತಕ್ಕೆ} \kn{ಅಜ್ಜಿ} \kn{ಶುಂಠಿಯ} \kn{ಕಷಾಯ} \kn{ಮಾಡಿದರು}.
\end{quote}

\begin{quote}
\emph{śītakke ajji śuṇṭhiya kaṣāya māḍidaru.}
\end{quote}

\begin{quote}
Grandmother made a ginger decoction for the cold.
\end{quote}

Boiled from ginger, pepper or tulasi leaves and drunk hot.

\begin{cousinweb}
\textbf{\kn{ಕಷಾಯ}} (\emph{kaṣāya}) is borrowed from Sanskrit \emph{kaṣāya}, \textquotedblleft{}astringent; an astringent juice\textquotedblright{} (CDIAL 2974).
\end{cousinweb}

\begin{grammarlens}[title={What you've built}]
One more word for describing what you lived and felt.
\end{grammarlens}

\subsection*{Your turn}
\begin{itemize}
\raggedright
  \item \emph{Say it:} \emph{kaṣāya}
  \item \emph{Say it:} \emph{kaṣāya}, once more
  \item \emph{Say it:} \emph{śītakke ajji śuṇṭhiya kaṣāya māḍidaru.} — Grandmother made a ginger decoction for the cold.
  \item \emph{Recall:} say the Kannada for a heartbeat, then the Kannada for nursing, then say \emph{kaṣāya} again
\end{itemize}

\begin{tcolorbox}[breakable,colback=teal!4,colframe=teal!35!black,title={Before you move on}]
What does \kn{ಕಷಾಯ} mean? (\textquotedblleft{}A herbal decoction\textquotedblright{}.) Say it once more.
\end{tcolorbox}
